\chapter{Analysis and Requirements Engineering}

\section{Stakeholder Analysis}

Primary stakeholders include visitors, curators, research engineers, and operations
teams. Each requires a different quality guarantee: usability, reliability,
explainability, and maintainability. Visitors prioritize speed and clarity; curators
emphasize factual integrity and update control; engineers prioritize maintainability
and reproducibility; leadership focuses on adoption and measurable engagement. A
system that optimizes only one of these viewpoints tends to fail in operation; hence
requirements engineering here explicitly balances all four.

\section{Functional Requirements}

\begin{enumerate}[label=\textbf{FR-\arabic*.},leftmargin=3em]
  \item \textbf{Real-time detection.} Detect and track artifacts in the live camera
        feed with bounding-box overlays at interactive frame rates.
  \item \textbf{On-device classification.} Classify a framed artifact as one of the
        51 catalogued GEM artifacts entirely on-device, with no per-frame network call.
  \item \textbf{Confidence gating.} A confident detection opens the artifact's
        information; a view that falls below the per-class confidence threshold is
        suppressed, so the app keeps scanning instead of guessing.
  \item \textbf{Artifact information.} Present name, period/dynasty, material, and
        description, with an original$\rightarrow$restored image toggle where available.
  \item \textbf{Bilingual AI guide.} Answer artifact-specific and general questions in
        the user's language, grounded in curated museum data.
  \item \textbf{Bilingual UI.} Support full English/Arabic switching including
        right-to-left layout, with the chosen language persisted across sessions.
  \item \textbf{Text-to-speech.} Optionally read AI replies aloud in the active language.
  \item \textbf{Photo capture.} Capture and save camera photos to the device gallery.
  \item \textbf{Museum information and location.} Provide a museum-information screen
        (hours, location, ticketing, map) with location gating for on-site versus
        off-site visitors.
\end{enumerate}

\section{Non-Functional Requirements}

\begin{enumerate}[label=\textbf{NFR-\arabic*.},leftmargin=3.4em]
  \item \textbf{Latency:} real-time on-device detection (target $\geq 2.5$\,FPS
        on a mid-tier Android phone) and AI answers within a few seconds.
  \item \textbf{Footprint:} compact TensorFlow Lite model ($\leq 8$\,MB; actual
        5.97\,MB), installable on low-storage devices.
  \item \textbf{Robustness:} tolerate clutter, overlap, oblique angles, and
        fine-grained confusion among near-identical artifacts.
  \item \textbf{Privacy/edge:} perform recognition on-device with no per-frame
        upload; serve the AI backend locally, avoiding any third-party LLM API.
  \item \textbf{Accessibility:} bilingual content plus text-to-speech audio.
  \item \textbf{Maintainability/scalability:} modular, independently replaceable
        services; the knowledge base is updatable without an app release.
\end{enumerate}

\section{Data Requirements}

\subsection{Data Sources and Acquisition}
Training data consists of roughly 30{,}000 frames extracted (at about one frame per
second) from approximately 50 in-gallery videos of GEM artifacts captured during
the museum's opening hours, spanning all 51 classes. This source reflects the real
deployment domain far better than generic web imagery.

\subsection{Volumetric and Diversity Requirements}
The dataset must cover class diversity, view diversity, and clutter diversity, with
balanced per-class representation. Per-class frame counts range from a few hundred to
just under a thousand, which is reasonably balanced for the task.

\subsection{Annotation and Labeling Standards}
Annotation consistency, class-map governance, and \emph{leakage-aware} split policies
are required. Because frames are extracted from continuous video, adjacent frames are
near-duplicates; splits must therefore be performed \emph{chronologically within each
class/video group} rather than by random shuffle, to avoid leaking near-identical
frames across train, validation, and test (see Chapter~5).

\section{App Persona}

To keep design decisions anchored to a concrete user rather than an abstract average,
we model a representative visitor (Figure~\ref{fig:persona}). Her goals and frustrations
motivate the system's emphasis on speed, offline robustness, and a low-friction
scan-to-answer path.

\begin{figure}[htbp]
\centering
\fbox{\begin{minipage}{0.9\textwidth}\small
\vspace{1mm}
{\large\bfseries Persona: Sarah the Curious Visitor}\\[1mm]
\textit{32 years old, international tourist, tech-savvy but time-constrained.}\\[2mm]
\begin{minipage}[t]{0.46\textwidth}
\textbf{Goals:}
\begin{itemize}[leftmargin=1.2em,nosep]
  \item Quickly identify artifacts.
  \item Learn historical context.
  \item Interactive discovery.
\end{itemize}
\end{minipage}\hfill
\begin{minipage}[t]{0.46\textwidth}
\textbf{Frustrations:}
\begin{itemize}[leftmargin=1.2em,nosep]
  \item Weak Wi-Fi in halls.
  \item Static museum text.
  \item High app latency.
\end{itemize}
\end{minipage}\\[2mm]
\textit{``I want to point my phone at something and instantly hear its story, like a
personal digital curator.''}
\vspace{1mm}
\end{minipage}}
\caption{User persona modeling for the \TimeLens{} experience design.}
\label{fig:persona}
\end{figure}

\section{Use Case Diagram}

The visitor is the primary actor, and the system exposes three core use cases that map
directly to the functional requirements (Figure~\ref{fig:usecase}): scanning to
identify an artifact (FR-01--FR-03), viewing its information (FR-04), and asking the
bilingual AI guide a question (FR-05). All three are reached from a single camera-first
screen, keeping the interaction model deliberately small.

\begin{figure}[htbp]
\centering
\begin{tikzpicture}[node distance=6mm and 28mm]
  % actor
  \node (head) at (0,0) [circle,draw,minimum size=5mm] {};
  \draw (0,-2.5mm)--(0,-12mm);            % body
  \draw (-5mm,-6mm)--(5mm,-6mm);          % arms
  \draw (0,-12mm)--(-4mm,-20mm);          % legs
  \draw (0,-12mm)--(4mm,-20mm);
  \node (visitor) at (0,-2.5) {\small Visitor};
  % use cases
  \node[draw,ellipse,align=center,minimum width=34mm,font=\small] (scan) at (4.2,1.2) {Scan / Identify\\ Artifact};
  \node[draw,ellipse,align=center,minimum width=34mm,font=\small] (info) at (4.2,-0.3) {View Artifact\\ Information};
  \node[draw,ellipse,align=center,minimum width=34mm,font=\small] (chat) at (4.2,-1.8) {Ask the AI Guide};
  % system boundary
  \begin{scope}[on background layer]
    \node[draw,dashed,rounded corners,fit=(scan)(info)(chat),inner sep=5mm,label=above:{\small TimeLens Platform}] {};
  \end{scope}
  % associations
  \draw (head) -- (scan);
  \draw (head) -- (info);
  \draw (head) -- (chat);
\end{tikzpicture}
\caption{Primary use cases for visitor interaction.}
\label{fig:usecase}
\end{figure}

\FloatBarrier

\section{Interaction and Data Modeling}

System interactions are organized by actor role. Data modeling focuses on the
alignment of user sessions, artifact scans, and bilingual content, ensuring
consistency between the on-device layer and the knowledge backend.

\section{Requirements Traceability Matrix}

Table~\ref{tab:traceability} maps a representative subset of requirements to the
components that satisfy them and the method used to validate each.

\begin{table}[htbp]
\centering\small
\caption{Requirement-to-component mapping.}
\label{tab:traceability}
\begin{tabularx}{\textwidth}{@{}l>{\raggedright\arraybackslash}X>{\raggedright\arraybackslash}X>{\raggedright\arraybackslash}p{2.3cm}@{}}
\toprule
\textbf{Req.} & \textbf{Description} & \textbf{Component} & \textbf{Validation}\\
\midrule
FR-01 & Live detection $+$ tracking overlay & ARService $+$ TFLiteService (YOLOv8n) & Functional test\\
FR-02 & On-device 51-class classification & TFLiteService (end-to-end NMS) & Per-class mAP eval\\
FR-03 & Confident $\rightarrow$ info; sub-threshold $\rightarrow$ suppressed & Home / Scanning logic & End-to-end scenario\\
FR-05 & Bilingual grounded AI guide & RAG \texttt{/ask} (ChromaDB $+$ Gemma) $+$ ChatDrawer & API integration\\
FR-06 & EN/AR switch $+$ RTL $+$ persistence & LocaleController $+$ l10n & UI test (both locales)\\
FR-07 & Text-to-speech replies & \texttt{flutter\_tts} (ChatDrawer) & Functional test\\
NFR-2 & Compact model $\leq 8$\,MB & YOLOv8n TFLite FP16 (5.97\,MB) & Asset-size check\\
NFR-4 & On-device privacy / local backend & On-device inference $+$ local Ollama & Architecture review\\
\bottomrule
\end{tabularx}
\end{table}
